\section{Introduction}\label{sec:intro}
% Draft 2026-10-17 (writer). Claims are listed in paper/claims.md (rows I1-I9). No results are stated as facts.

An in-memory key--value store such as Redis is often the first component to feel a change in load. When it runs
on Kubernetes, the operator chooses how many pods serve the cluster. Too few pods and requests queue up, so tail
latency grows and the service misses its latency objective~\cite{dean2013tail}. Too many pods and the operator
pays for capacity that sits idle. The default answer is the Horizontal Pod Autoscaler (HPA). It runs a control
loop, by default every 15 seconds, and sets the replica count from the ratio between the current value of a
metric and its target~\cite{k8sdocs_hpa}. The rule is simple, but it can only respond to load that has
already arrived. A pod added now serves traffic only after it is scheduled, its image is pulled and it becomes
ready~\cite{k8sslos,harter2016slacker}, and in Redis Cluster hash slots must then be moved to it before it takes
its share of the keys~\cite{redisdocs_scaling,redisdocs_clusterspec}.

Learning-based autoscalers try to act before the load arrives. DERP~\cite{bitsakos2018derp} applied deep
reinforcement learning to elastic resource provisioning. DInos~\cite{bitsakos2025dinos} targets autoscaling
that generalizes across stateless cloud applications. Its agent reads a window of recent metrics through a Long Short-Term Memory (LSTM)
network~\cite{hochreiter1997lstm}, is trained in a simulator, and is then fine-tuned for a new environment,
including a Kubernetes cluster running Redis. That design has a cost that is easy to overlook. In the DInos
implementation the Q-network has \resParamsDInos{} trainable parameters, and almost all of them belong to the
recurrent encoder. Fine-tuning for a new environment updates every one of them.

Variational quantum circuits (VQCs) have been used as function approximators in reinforcement learning, both as
policies~\cite{jerbi2021parametrized} and as Q-functions~\cite{skolik2022quantum}, and hybrid networks that put a
trainable circuit after a pre-trained classical network have been used for transfer
learning~\cite{mari2020transfer}. The circuits in these studies are small. A small parameter count alone shows
little, though. Skolik et al.\ report that architectural and hyperparameter choices matter more than
the number of parameters~\cite{skolik2022quantum}, so a small classical network could match a small circuit. We
therefore treat parameter efficiency as a hypothesis and test it against a classical control with the same number
of trainable parameters. We also keep timing claims apart from it. Our circuits run on a classical simulator, and
simulator wall-clock time says nothing about speed on quantum hardware.

A comparison of this kind is only as good as the simulator behind it. If latency or CPU did not depend on the
number of pods, the cheapest policy would be to keep one pod at all times, and a tiny model of any kind would
look as good as a large one. We therefore built a new simulator for this study. Arrivals enter a queue served by
the ready pods, so latency and CPU respond to the scaling action. New pods start after a delay, and capacity drops
for a while after each change while the cluster rebalances. Before any agent is compared, the simulator must pass
fixed sanity gates: always using the minimum or the maximum number of pods, or the best static pod count, has to
lose to a tuned threshold policy, and an oracle that knows future arrivals gives an upper bound on the
reward.\todo{D8: the researcher decides whether and how the paper states why the earlier simulator was replaced.}

This paper presents Sofia, a hybrid quantum--classical deep reinforcement learning agent for Kubernetes
autoscaling. Sofia keeps the DInos interface: the same window of cluster metrics as input and the same three
actions (add a pod, remove a pod, keep the current count). It replaces part of the DInos Q-network with a VQC
of \resQubits{} qubits.\todo{Which part (D5): full Q-function, head after a small encoder, or the fine-tuned
adaptation layer. Rewrite this sentence and C2 below once decided.}

The paper makes four contributions. First, we describe a simulator for Redis on Kubernetes in which the reward
depends on the scaling action through queueing, start-up delay and rebalancing, together with the sanity gates and
the oracle bound that any agent comparison must pass (Section~\ref{sec:simv2}). Second, we design Sofia and
specify its circuit, its encoding of the metric window and its training procedure
(Section~\ref{sec:quantum}). Third, we compare Sofia with DInos retrained in the same simulator, a threshold
policy modelled on the HPA, and a parameter-matched classical control, on a periodic workload with noise and
spikes as in DInos and on a bursty Poisson-based workload\todo{Workload list per D3 (W0--W3).}, over
\resNumSeeds{} seeds (Section~\ref{sec:eval}). We report SLA violations, p95 latency, steps to reach a target
reward, trainable parameters and checkpoint size. Fourth, we check whether the ranking of the agents carries over
to a Kubernetes cluster running Redis.\todo{Keep only if real-cluster runs are in this paper (PLAN 8.7, G6).}
Our circuits run on a classical simulator\todo{Add IBM hardware agreement (\resHWAgreement) only if D9 runs are
done.}, and none of our results is a claim about quantum speed-up.

Section~\ref{sec:bg} gives background on Kubernetes autoscaling, Redis and variational circuits.
Section~\ref{sec:related} reviews related work. Section~\ref{sec:system} describes the system and
Section~\ref{sec:simv2} the simulator. Sections~\ref{sec:agents} and~\ref{sec:quantum} present the agents,
the baselines and the quantum component. Section~\ref{sec:eval} reports the evaluation,
Section~\ref{sec:disc} discusses limitations and Section~\ref{sec:concl} concludes.
